\section{Data and Evaluation Protocol}
Table~\ref{t2} lists the six main contact streams and the discussion archives. Dept1, used for the contact-selectivity diagnostic, comes from the same institutional email data as Dept2 and Dept4~\cite{motifs}. Company and Enron record organizational email~\cite{company,enron}, and CollegeMsg records private messages~\cite{college}. GNU libc supplies both contacts and threads~\cite{libc}. The other discussions come from Python-Dev and Python-Ideas~\cite{pythondev,pythonideas}, R-devel and R-package-devel~\cite{rdevel,rpackagedevel}, and Samba~\cite{sambatechnical}.
\resulttable{tab2_data}{Network and thread populations. Above: distinct nodes and cleaned contacts (GNU libc: all resolved 2021--22 replies). Below: posting addresses and messages within the evaluation period. Windows include scheduled windows or months with no eligible events.}{t2}{4pt}

Each 14-day simulation window uses contacts within its initial roster. GNU libc uses a rolling 98-day roster at each month boundary. Each development network supplies 1,000 training and 300 subsequent validation events. We select among epochs 3, 7 and 12 by true-pair log loss (AdamW, learning rate 0.002, weight decay 0.005, batch size 24). Dept2 and Company use local weights. All other networks use frozen Dept2 weights. Only the original Dept4 comparison fixed these weights and windows before target-label generation. The Enron reevaluation and feature ablations use previously inspected networks.

At budget \(K\), joint and at-least-one coverage average \(\mathbf{1}\{S\subseteq L\}\) and \(\mathbf{1}\{S\cap L\ne\varnothing\}\). Pair MRR averages reciprocal exact-pair rank. Per-source recall is \(|S\cap L|/k\). Five seeds share one event denominator. Paired 95\% intervals use 10,000 bootstrap replicates over seeds and circular three-window blocks, including empty windows, with event-weighted means. Four windows yield little temporal variation. With three Python follow-up months, each block covers all months, so only seed variability contributes.

\section{Retrieval on Simulated Cascades}
\subsection{What a fixed budget buys}
Across six networks, 8,498 uniformly seeded cascades qualified (Table~\ref{t3}). The average CollegeMsg pool was 1,830, against 115 in Dept4. At 12 places the CollegeMsg MLP recovered any source in 6.88\% of events and both in none. Even 24 places gave only 0.08\% joint coverage. Budgets of 5\% and 10\% raised coverage to 2.21\% and 5.21\%. In every other network 5\% afforded less coverage than 12 places (\figref{f2}).
\resulttable{tab3_main}{Retrieval by network. Pool is the mean legal roster. Any@12 (full MLP) and Both columns are percentages. Static is undirected and additive. Tree uses full features. Stars indicate one archived seed, dashes unavailable comparisons.}{t3}{3.3pt}
\paperfig{fig2_budgets}{Joint coverage at fixed and proportional budgets (rounded up per event). Points are full-MLP seed means with 95\% seed-and-window intervals.}{f2}

On the prespecified Dept4 transfer, seven of 12 windows yielded 1,308 events. Full MLP coverage reached 23.49\%, versus 20.87\% for the static model and 17.87\% for five-frame FGSSI. Paired gains were 2.61 points [1.66, 3.89] and 5.61 [4.44, 6.98]. Although the ranking transferred between departments, exact-pair MRR was only 0.03055.

\subsection{When do ties to visible participants distinguish sources?}
In Dept2, 52.2\% of sources, against 8.6\% of alternatives, had sent a message to a visible member in the preceding 28 days. Adding contact features raised joint coverage from 1.28\% to 14.68\% (+13.40 points [11.42, 15.48]). In Company such ties were common among both sources (67.7\%) and alternatives (55.1\%). They yielded only 0.02 points [$-0.78$, 0.75].

Figure~\ref{f3} compares outgoing contact counts within each event (ties earn half a win). Mean AUC was 0.719 on Dept2 and 0.556 on Company. Across seven simulated networks, including Dept1, AUC correlated positively with gain over personal history (Pearson~\(r=0.818\)). Across seven real list periods, its correlation with gain over posting volume was strongly negative (\(r=-0.901\)). GNU libc 2021--22 had the highest AUC (0.754), yet frequent posters were so often the early speakers that volume alone retrieved both in 40.23\% of threads, against 19.20\% for the simulation-trained contact model. A concentration of activity among core maintainers could explain this pattern, although the archive records posting activity rather than maintainer status. Using personal history as the baseline throughout gave~\(r=0.836\) across the 14 domains. Related periods and networks contribute dependent points.
\paperfig{fig3_selectivity}{Within-event AUC of 28-day outgoing contacts versus gain over each baseline. Each network or period has equal weight. Real AUC requires two reachable targets, whereas coverage counts all eligible threads.}{f3}

Histogram gradient boosting on the same masks yielded contact-over-personal gains of 12.62 points [10.65, 14.67] on Dept2 and 0.13 [$-0.68$, 0.87] on Company. The contrast survived a change of training objective: the trees fit candidate-level binary labels using all legal negatives, with 24 or 48 rounds selected by validation pair loss.

Table~\ref{t4} varies propagation on coupled draws. Sources remained active, recovered after 48 or 72 hours, or spread alongside spontaneous entries (2\% per nonsource per window). Each rule had 100 attempts per window across four Dept2 and eight Enron windows. On draws qualifying under all rules, all four Dept2 gains remained positive. All four Enron intervals included zero. Rule-specific cohorts also reflect changes in which draws produce enough observable activity.
\resulttable{tab4_rules}{Contact-minus-personal joint coverage (percentage points). All uses each rule's eligible draws. Joint uses their intersection. Intervals are paired within each cohort.}{t4}{3pt}

\subsection{Source priors, silent partners, and source count}\label{sec:prior}
Let \(w_i\) be one plus candidate \(i\)'s outgoing count in the preceding 28 days. Assign distinct pairs weights \(f_0=1\), \(f_1=w_i+w_j\), and \(f_2=w_iw_j\). For component \(r\), let \(\bar f_r\) be its mean weight over all roster pairs before masking. At mixture weights \(\alpha_r\), an explicit offset multiplies the additive pair mass by \(\sum_r\alpha_r f_r/\bar f_r\). The direct model and explicit model see the same mixed training and validation events. The latter fixes \(\alpha_r=1/3\) in training and inference. A deployment-only control trains on uniform events and applies the same equal-mixture offset at prediction time.

Under Enron test weights \((0.05,0.15,0.80)\), direct mixed training retrieved both sources in 34.05\% of 2,607 events. Explicit training reached 39.75\%, a paired gain of 5.70 points [2.88, 9.39]. Deployment-only correction reached 38.93\%, leaving a smaller explicit-training advantage of 0.83 points [$-0.65$, 2.39]. When source pairs were instead drawn preferentially along a recent direct tie, the offset omitted the dependence driving selection and fell 2.24 points below direct training [$-3.31$, $-1.02$]. Figure~\ref{f4} contrasts these empirical results with expected coverage conditional on a complete trace.
\paperfig{fig4_prior_and_ceiling}{Left, Enron coverage across four source laws with marginal 95\% seed-and-window intervals. Right, one-emitter complete-trace ceilings under the product prior, partitioned into captured mass, emitter omission and partner-ranking loss. D2, Co., D4 and Enr. denote Dept2, Company, Dept4 and Enron, respectively.}{f4}

On Enron the mean oracle ceiling for a silent partner was 14.22\% under uniform draws (778 one-emitter events) and 51.73\% under the pure product law (1,102 events). In the latter cohort a suffix model calibrated on 32 earlier labeled Dept2 events retained 19.36\% conditional coverage. Missing the emitter cost 7.46 points and partner ranking cost 24.90. The populations differ across laws because the pair is drawn before the cascade qualifies. On Company, the product-prior ceiling was 88.51\% while the public model captured only 6.33\%. Given the full contact stream and exact join times, we marginalize over parent identities: a susceptible recipient joins a batch with probability \(1-(1-p)^m\) after \(m\) contacts from active senders. We evaluate the resulting list under the full-trace posterior. Revealing successful parents increased optimal list coverage by more than \(10^{-9}\) on only 3 of 12,280 trajectories, all on Company.

With three uniform sources, per-source recall reached 39.23\% on Dept2 and 28.46\% on Enron, but all-source coverage was only 1.30\% and 0.74\%. In a separate arm we drew \(k\) uniformly from one to three per attempt and withheld it from the frozen scorer. Figure~\ref{f5} stratifies the accepted events by their actual count. For Enron, all-source coverage falls from 68.06\% at \(k=1\) to 4.64\% at \(k=2\) and 1.08\% at \(k=3\). Per-source recall falls more slowly. A single-source eligible cascade necessarily contains a successful transmission, whereas nearly nine in ten eligible three-source events had a silent member.
\paperfig{fig5_multisource}{Source count hidden from the ranker. Left, fraction of eligible events with a silent source. Right, full-model all-source coverage and per-source recall, equal at \(k=1\).}{f5}

Two exploratory tests probe remaining ambiguities. Estimating the prior mixture by EM from up to 32 unlabeled suffixes per window changed coverage by \(+1.03\) percentage points on Dept4 and \(-4.41\) percentage points on Enron relative to the equal mixture under the 80\% two-active law. For an exact-pair decision, a single-snapshot adaptation of BOSouL and a learned snapshot model each had zero exact hits on 100 Dept2 and 100 Company events. The search used 20 initial evaluations, 10 search steps and 10 simulations per evaluation. These cohorts support neither reliable unlabeled adaptation across networks nor an exact-pair advantage at that search budget.

\section{Transfer to Observed Discussions}
\subsection{D0 retrieval across five held-out community periods}
At D0 we hide the root author and first distinct respondent and retain all later first-appearance batches. Table~\ref{t5} evaluates the role models of Sect.~3.3 on 1,793 threads. Scoring requires one visible roster member and 12 legal candidates. Out-of-roster targets and abstentions count as misses. The coefficients were learned on earlier Python threads. The original R and GNU libc slices had been inspected previously. GNU libc 2024 and Samba 2020 test transfer to a later period and another community, respectively.
\resulttable{tab5_heldout}{D0 threads with both speakers retrieved. Reachable requires a scorable suffix and both targets in the pool. The four columns rank by 28-day posts, historical roles, roles plus replies, and the full enrichment model.}{t5}{4pt}

On GNU libc, roles retrieved 375 versus 354 pairs for volume in 2021--22, and 136 versus 122 in 2024. Adding reply and enrichment features yielded 383 and 134 hits. On Samba the full model retrieved 22 pairs, compared with 14 for volume and 23 for direct replies. Its gain over volume was 10.13 points [$-1.12$, 27.42]. Richer evidence about later speakers did not consistently improve on the role history.

\subsection{Masking the same discussions more deeply}
D25 and D50 hide additional whole batches containing at most one quarter or one half of later joiners. Positive depths require one hidden batch and three remaining participants. Table~\ref{t6} and \figref{f6} follow the same eligible threads at all three depths. Every depth now requires three visible roster members to score, and the full model is the frozen simulation-trained MLP. Both the D0 denominator and the full model therefore differ from Table~\ref{t5}.
\resulttable{tab6_depth}{Full temporal MLP minus post volume in percentage points on identical threads at D0, D25 and D50. A thread may become unscorable after masking and stays in the denominator.}{t6}{7pt}
\paperfig{fig6_real_depth}{Joint coverage (\%) on the fixed threads of Table~\ref{t6}. Coincident markers indicate ties. Each list period has its own vertical scale.}{f6}

On the same 24 Python-Dev threads, volume recovered five pairs at both D25 and D50. MLP coverage fell from 25.83\% to 11.67\%. Its paired gain changed from \(+5.00\) points [2.50, 7.50] to \(-9.17\) [$-15.83$, $-4.17$]. The reversal survives fixing thread identity. Masking can still remove roster members: at D50, nine of 17 R-devel threads were scorable and none had both targets in the pool. Every method consequently scored zero.

\subsection{A controlled change in the role of the second target}
We crossed a uniform or activity-weighted first sender with two partner definitions. In the symmetric condition both sources start active. In the asymmetric condition the sender starts alone and its unique first activation becomes the partner. Holding the scorer fixed, we retain draws eligible in all four conditions and average each intervention over its two possible orders. The two resulting terms divide the interaction equally.

In both networks at both depths, requiring the partner to participate in a successful contact mattered more than weighting the first sender by activity (\figref{f7}). At D50 the role term was 55.55 points on Dept2 (110 common draws) and 29.05 on Enron (42), versus activity terms of 7.36 and 2.86. This decomposition measures interventions in the simulator. The observed archives do not identify their contributions to the real retrieval gap.
\paperfig{fig7_sim2real_decomp}{Simulated effects of the sender prior and partner role on joint coverage, averaged over both intervention orders on common draws (paired 95\% intervals).}{f7}

\section{Discussion and Conclusion}
\paragraph{Limitations.}
The simulator selects sources from the historical roster and conditions evaluation on enough subsequent activity. Changing the number of sources or their prior can alter both evidence and eligibility. The complete-trace posterior is exact for independent, constant-probability synchronous attempts under the stated model. Correlated contacts or individual transmission probabilities would change its failure-count likelihood. The departmental networks share an institution and the mailing-list periods share contributors. Network-level correlations and short-window bootstrap intervals have limited support for extrapolation to another organization.

\paragraph{Ethical and Privacy Considerations.}
Historical correspondence reveals social relationships even when message bodies are excluded. We use header-derived reply edges and hashed message and address identifiers. Hashing can preserve linkability. We therefore report aggregate coverage and omit individual candidate lists, rare thread examples and address-level predictions. A root author and an early respondent are observable roles. Assigning responsibility for an idea or an event would require independent evidence. An address can also split one person across aliases or conceal activity outside the archive.

\paragraph{Conclusion.}
A late suffix often identifies one emitter while leaving many partners plausible. The complete-trace posterior quantifies the ambiguity that survives even observing every success. Historical ties improve shortlists where those ties are rare among competing candidates. In real discussions, early speaking roles and posting volume provide different evidence, and masking can erase the advantage of a simulation-trained ranking even on the same threads.
